As specified in the previous paragraph \ref{subsub:implementationMpcioAndCommunication}, \code{prac} is executed three times simultaneously.
Hence, assuming one executes locally without using docker, one needs to run \code{prac} in three terminals at once.
Calling \code{prac} requires several arguments: player number, which mode to execute (-p, -o, or blank), number of threads for faster execution (-t), and for the server $P_2$ the type, amount and depth of functions to calculate.
\code{r} will represent RDPFs (\ref{sub:RDPFTriples}).

Thus, the function calls for preprocessing will look as follows:
\begin{enumerate}
    \item \code{./prac -t 3 -p 0}
    \item \code{./prac -t 3 -p 1 localhost}
    \item \code{./prac -t 3 -p 2 localhost localhost r20:20}
\end{enumerate}
This will execute the precomputation generating 20 RDPF-triples, each of depth 20, with each instance working in three threads.

\begin{enumerate}
    \item \code{./prac -t 3 0 duoram 20 5}
    \item \code{./prac -t 3 1 localhost duoram 20 5}
    \item \code{./prac -t 3 2 localhost localhost duoram 20 5}
\end{enumerate}
Will execute the online phase, again with three threads each.
Specifying \code{duoram} will execute the tests as described in \ref{sub:testingAndBenchmarking} with 5 messages using RDPFs of depth 20 for each subtest.

\subsubsection*{General Execution}
All executions start in \code{prac.cpp} which accomodates the \code{main} method.
This includes the starting points for both the players' and the server's computations.
If parses the given arguments, sets the adresses of the other parties, and deceides whether the current execution is a player's or server's computation.
That means it calls \code{comp\_player\_main} or \code{server\_player\_main}.

\code{comp\_player\_main} represents the player's computation.
It is, again, separated into \code{preprocessing\_comp} and \code{online\_main} representing the preprocessing computation and the computations necessary for performing the onlien phase.

\code{server\_player\_main} represents the servers' computations.
It too is separated into its phases \code{preprocessing\_server} and \code{online\_main}.

The careful reader notices both the server and the player calling \code{comp\_main} albeit the previous text assumed these computations to be implemented separately.
Notice that this too is only a starting point for more specialized calculations.
Thus, \code{online\_main} too further specifies its computations given the number of the computational peer.

\subsubsection*{Preprocessing} \label{subsub:implementationPreprocessing}
%TODO wie wird das in files gespeichert und wo findet man die?


\subsubsection*{Cyclic Shifting} \label{subsub:implementationCyclicShifting}
\subsubsection*{Online Phase} \label{subsub:implementationOnlinePhase}
%TODO wie werden die DPFs (RDPF, CDPF) aufgerufen
\subsubsection*{Reading} \label{subsub:implementationReading}
\subsubsection*{Writing} \label{subsub:implementationWriting}